% GENERATED FILE. Edit canonical lessons, then run npm run generate:books.
% canonical-source-hash: fnv1a64:ac289c9240d082e7
% canonical-lessons: JA-C100-gomi, JA-C100-bou, JA-C100-ita, JA-C100-hasami, JA-C100-zasshi

\chapter{Rubbish, Stick, Board, Scissors, Magazine}
\label{ch:ja-rubbish-stick-board-scissors-magazine}
\hlchaptermodality{100}

\begin{chapteropening}
\textbf{By the end of this chapter:} \emph{I can say rubbish, a stick, a board, scissors and a magazine.}

Complete the last lesson of chapter 100: i can say rubbish, a stick, a board, scissors and a magazine.
\end{chapteropening}

\section[\texorpdfstring{\ja{ごみ} (gomi)}{ごみ (gomi)}]{\ja{ごみ} (gomi) — rubbish}

\label{lesson:JA-C100-gomi}



\begin{quote}
Before the new one: say the Japanese for a scrubbing brush, then the Japanese for a dustpan.
\end{quote}

\subsection*{You'll want to know: \ja{ごみ}}
\textbf{\ja{ごみ}} — \emph{gomi} — \textquotedblleft{}rubbish\textquotedblright{}.

\begin{grammarlens}[title={What you've built}]
One more word for this chapter.
\end{grammarlens}

\subsection*{Your turn}
\begin{itemize}
\raggedright
  \item \emph{Say it:} \emph{gomi}
  \item \emph{Say it:} \emph{gomi}, once more
  \item \emph{Say it:} say \emph{gomi}
  \item \emph{Recall:} say the Japanese for a scrubbing brush, then the Japanese for a dustpan, then say \emph{gomi} again
\end{itemize}

\begin{tcolorbox}[breakable,colback=teal!4,colframe=teal!35!black,title={Before you move on}]
What does \ja{ごみ} mean? (\textquotedblleft{}Rubbish\textquotedblright{}.) Say it once more.
\end{tcolorbox}
\section[\texorpdfstring{\ja{ぼう} (bou)}{ぼう (bou)}]{\ja{ぼう} (bou) — a stick, a pole}

\label{lesson:JA-C100-bou}



\begin{quote}
Before the new one: say the Japanese for a dustpan, then the Japanese for rubbish.
\end{quote}

\subsection*{You'll want to know: \ja{ぼう}}
\textbf{\ja{ぼう}} — \emph{bou} — \textquotedblleft{}a stick, a pole\textquotedblright{}.

\begin{grammarlens}[title={What you've built}]
One more word for this chapter.
\end{grammarlens}

\subsection*{Your turn}
\begin{itemize}
\raggedright
  \item \emph{Say it:} \emph{bou}
  \item \emph{Say it:} \emph{bou}, once more
  \item \emph{Say it:} say \emph{bou}
  \item \emph{Recall:} say the Japanese for a dustpan, then the Japanese for rubbish, then say \emph{bou} again
\end{itemize}

\begin{tcolorbox}[breakable,colback=teal!4,colframe=teal!35!black,title={Before you move on}]
What does \ja{ぼう} mean? (\textquotedblleft{}A stick, a pole\textquotedblright{}.) Say it once more.
\end{tcolorbox}
\section[\texorpdfstring{\ja{いた} (ita)}{いた (ita)}]{\ja{いた} (ita) — a board}

\label{lesson:JA-C100-ita}



\begin{quote}
Before the new one: say the Japanese for rubbish, then the Japanese for a stick.
\end{quote}

\subsection*{You'll want to know: \ja{いた}}
\textbf{\ja{いた}} — \emph{ita} — \textquotedblleft{}a board\textquotedblright{}.

\begin{grammarlens}[title={What you've built}]
One more word for this chapter.
\end{grammarlens}

\subsection*{Your turn}
\begin{itemize}
\raggedright
  \item \emph{Say it:} \emph{ita}
  \item \emph{Say it:} \emph{ita}, once more
  \item \emph{Say it:} say \emph{ita}
  \item \emph{Recall:} say the Japanese for rubbish, then the Japanese for a stick, then say \emph{ita} again
\end{itemize}

\begin{tcolorbox}[breakable,colback=teal!4,colframe=teal!35!black,title={Before you move on}]
What does \ja{いた} mean? (\textquotedblleft{}A board\textquotedblright{}.) Say it once more.
\end{tcolorbox}
\section[\texorpdfstring{\ja{はさみ} (hasami)}{はさみ (hasami)}]{\ja{はさみ} (hasami) — scissors}

\label{lesson:JA-C100-hasami}



\begin{quote}
Before the new one: say the Japanese for a stick, then the Japanese for a board.
\end{quote}

\subsection*{You'll want to know: \ja{はさみ}}
\textbf{\ja{はさみ}} — \emph{hasami} — \textquotedblleft{}scissors\textquotedblright{}.

\begin{grammarlens}[title={What you've built}]
One more word for this chapter.
\end{grammarlens}

\subsection*{Your turn}
\begin{itemize}
\raggedright
  \item \emph{Say it:} \emph{hasami}
  \item \emph{Say it:} \emph{hasami}, once more
  \item \emph{Say it:} say \emph{hasami}
  \item \emph{Recall:} say the Japanese for a stick, then the Japanese for a board, then say \emph{hasami} again
\end{itemize}

\begin{tcolorbox}[breakable,colback=teal!4,colframe=teal!35!black,title={Before you move on}]
What does \ja{はさみ} mean? (\textquotedblleft{}Scissors\textquotedblright{}.) Say it once more.
\end{tcolorbox}
\section[\texorpdfstring{\ja{ざっし} (zasshi)}{ざっし (zasshi)}]{\ja{ざっし} (zasshi) — a magazine}

\label{lesson:JA-C100-zasshi}



\begin{quote}
Before the new one: say the Japanese for a board, then the Japanese for scissors.
\end{quote}

\subsection*{You'll want to know: \ja{ざっし}}
\textbf{\ja{ざっし}} — \emph{zasshi} — \textquotedblleft{}a magazine\textquotedblright{}.

\begin{grammarlens}[title={What you've built}]
One more word for this chapter.
\end{grammarlens}

\subsection*{Your turn}
\begin{itemize}
\raggedright
  \item \emph{Say it:} \emph{zasshi}
  \item \emph{Say it:} \emph{zasshi}, once more
  \item \emph{Say it:} say \emph{zasshi}
  \item \emph{Recall:} say the Japanese for a board, then the Japanese for scissors, then say \emph{zasshi} again
\end{itemize}

\begin{tcolorbox}[breakable,colback=teal!4,colframe=teal!35!black,title={Before you move on}]
What does \ja{ざっし} mean? (\textquotedblleft{}A magazine\textquotedblright{}.) Say it once more.
\end{tcolorbox}
